\section{fmm3d\_\-mpi/tt0.cpp File Reference}
\label{tt0_8cpp}\index{fmm3d_mpi/tt0.cpp@{fmm3d\_\-mpi/tt0.cpp}}


{\tt \#include \char`\"{}fmm3d\_\-mpi.hpp\char`\"{}}\par
\subsection*{Namespaces}
\begin{CompactItemize}
\item 
namespace {\bf std}
\end{CompactItemize}
\subsection*{Functions}
\begin{CompactItemize}
\item 
int {\bf main} (int argc, char $\ast$$\ast$argv)
\end{CompactItemize}


\subsection{Function Documentation}
\index{tt0.cpp@{tt0.cpp}!main@{main}}
\index{main@{main}!tt0.cpp@{tt0.cpp}}
\subsubsection{\setlength{\rightskip}{0pt plus 5cm}int main (int {\em argc}, char $\ast$$\ast$ {\em argv})}\label{tt0_8cpp_30c4c19ff69d905e8e4c750b3af23485_30c4c19ff69d905e8e4c750b3af23485}




Definition at line 22 of file tt0.cpp.

References p\-A, and p\-C.

\footnotesize\begin{verbatim}23 {
24   MPI_Init( &argc, &argv );
25   PetscInitialize(&argc,&argv,"options_0",NULL);  //PetscTruth flg = PETSC_FALSE;
26   
27   MPI_Comm comm; comm = PETSC_COMM_WORLD;
28   int mpirank;  MPI_Comm_rank(MPI_COMM_WORLD, &mpirank);
29   int mpisize;  MPI_Comm_size(MPI_COMM_WORLD, &mpisize);
30   int dim = 3;
31   srand48( (long)( time(NULL)+mpirank ) );
32   
33   PetscTruth flg = PETSC_FALSE;
34   pC( PetscPrintf(MPI_COMM_WORLD, "mpisize %d\n", mpisize) );
35   printf( "Hello world from process %d of %d\n", mpirank, mpisize );
36   
37   //1. allocate random data
38   int numsrc;  pC( PetscOptionsGetInt(PETSC_NULL, "-numsrc", &numsrc, &flg) ); pA(flg==PETSC_TRUE);
39   int numtrg;  pC( PetscOptionsGetInt(PETSC_NULL, "-numtrg", &numtrg, &flg) ); pA(flg==PETSC_TRUE);
40   int kt;  pC( PetscOptionsGetInt(PETSC_NULL, "-kt", &kt, &flg) ); pA(flg==PETSC_TRUE);
41   
42   vector<double> tmp(2);         tmp[0] = 1;     tmp[1] = 0.25; //coefs in the kernel, work for all examples
43   Kernel3d_MPI knl(kt, tmp);
44   
45   int lclnumsrc = (numsrc+mpirank)/mpisize;
46   int lclnumtrg = (numtrg+mpirank)/mpisize;
47 
48   printf( "Hello world from process %d of %d\n", mpirank, lclnumsrc );
49   pC( PetscPrintf(MPI_COMM_WORLD, "luclnumsrc and lclnumtrg %d %d\n", lclnumsrc, lclnumtrg) );
50   
51   Vec srcPos;  pC( VecCreateMPI(comm, lclnumsrc*dim, PETSC_DETERMINE, &srcPos) );
52   double* srcPosarr; pC( VecGetArray(srcPos, &srcPosarr) );
53   for(int k=0; k<lclnumsrc*dim; k++){
54          srcPosarr[k] = (2.0*drand48()-1.0);
55   }
56   pC( VecRestoreArray(srcPos, &srcPosarr) );
57   
58   Vec trgPos;  pC( VecCreateMPI(comm, lclnumtrg*dim, PETSC_DETERMINE, &trgPos) );
59   double* trgPosarr; pC( VecGetArray(trgPos, &trgPosarr) );
60   for(int k=0; k<lclnumtrg*dim; k++)
61          trgPosarr[k] = (2.0*drand48()-1.0);
62   pC( VecRestoreArray(trgPos, &trgPosarr) );
63   
64   int srcDOF = knl.srcDOF();
65   int trgDOF = knl.trgDOF();
66   Vec srcDen;  pC( VecCreateMPI(comm, lclnumsrc*srcDOF, PETSC_DETERMINE, &srcDen) );
67   double* srcDenarr; pC( VecGetArray(srcDen, &srcDenarr) );
68   for(int k=0; k<lclnumsrc*srcDOF; k++)
69          srcDenarr[k] = drand48();
70   pC( VecRestoreArray(srcDen, &srcDenarr) );
71   
72   Vec trgVal;  pC( VecCreateMPI(comm, lclnumtrg*trgDOF, PETSC_DETERMINE, &trgVal) );
73   PetscScalar zero = 0;  pC( VecSet(trgVal, zero) );
74   
75   //2. allocate fmm 
76   clock_t ck0, ck1;
77   
78   FMM3d_MPI* fmm = new FMM3d_MPI("fmm3d_");
79   fmm->srcPos()=srcPos;
80   fmm->srcNor()=srcPos;
81   fmm->trgPos()=trgPos;
82   fmm->ctr() = Point3(0,0,0); // CENTER OF THE TOPLEVEL BOX
83   fmm->rootLevel() = 0;         // 2^(-rootlvl) is the RADIUS OF THE TOPLEVEL BOX
84   fmm->knl() = knl;
85   
86   ck0 = clock();
87   pC( fmm->setup() );
88   ck1 = clock();  
89   pC( PetscPrintf(MPI_COMM_WORLD, "fmm setup used %e secs\n", double(ck1-ck0)/CLOCKS_PER_SEC) );
90   
91   //3. run fmm
92   for(int i=0; i<3; i++) {
93          ck0 = clock();
94          pC( fmm->evaluate(srcDen, trgVal) );
95          ck1 = clock();
96          pC( PetscPrintf(MPI_COMM_WORLD, "fmm eval used %e secs\n", double(ck1-ck0)/CLOCKS_PER_SEC) );
97   }
98   
99   //4. check
100   ck0 = clock();
101   double rerr;
102   pC( fmm->check(srcDen, trgVal, 20, rerr) );
103   pC( PetscPrintf(MPI_COMM_WORLD, "relative %e\n", rerr) );
104   ck1 = clock();
105   pC( PetscPrintf(MPI_COMM_WORLD, "fmm check used %e secs\n", double(ck1-ck0)/CLOCKS_PER_SEC) );
106   
107   delete fmm;
108   int tmp2;
109   pC( VecGetSize(srcDen,&tmp2) );
110   pC( PetscPrintf(MPI_COMM_WORLD, "%d \n", tmp2 ));
111   pC( VecGetLocalSize(srcDen,&tmp2) );
112   pC( PetscPrintf(MPI_COMM_WORLD, "%d \n", tmp2 ));
113   pC( VecDestroy(srcPos) );
114   pC( VecDestroy(trgPos) );
115   pC( VecDestroy(srcDen) );
116   pC( VecDestroy(trgVal) );
117   
118   PetscFinalize();
119   return 0;
120 }
\end{verbatim}\normalsize 


